\documentclass[11pt]{article}
\usepackage[margin=1in]{geometry}
\usepackage{amsmath,amssymb,amsthm,booktabs,hyperref}
\usepackage[T1]{fontenc}
\usepackage{lmodern}
\hypersetup{colorlinks=true,linkcolor=blue,urlcolor=blue,citecolor=blue}

\newtheorem{theorem}{Theorem}
\newtheorem{proposition}[theorem]{Proposition}
\newtheorem{corollary}[theorem]{Corollary}
\newtheorem{remark}[theorem]{Remark}
\newcommand{\Q}{\mathbb Q}
\newcommand{\F}{\mathbb F}
\newcommand{\Sel}{\operatorname{Sel}}
\newcommand{\Sha}{\mathop{\mathrm{Sha}}}
\newcommand{\rank}{\operatorname{rank}}

\title{Three Independent Rational Sections in a Nonisotrivial Elliptic Family\\
and a Conditional 2-Selmer Density Criterion}
\author{ResearchOS / TECT-X working research note\\
\small Authorship, original-contribution attribution and specialist review pending}
\date{19 November 2026\\[0.2em]\small Preliminary draft --- not submitted or peer reviewed}

\begin{document}
\maketitle
\begin{abstract}
We exhibit a nonisotrivial family of elliptic curves over $\Q$ with three
explicit rational points whose images in the Mordell--Weil group modulo
doubling are independent for every member of a specified arithmetic
progression. In particular, the Mordell--Weil rank is at least three
throughout that progression. An additional rational point on its first
fiber gives an independently checkable rank-at-least-four certificate.
We then prove an elementary conditional arithmetic-statistics
statement: an average 2-Selmer size strictly smaller than $16$ along
the progression would force a positive proportion of the curves to
have 2-Selmer dimension and Mordell--Weil rank exactly three.
\emph{No such average bound, global Selmer exhaustion theorem,
analytic-rank statement or Birch--Swinnerton-Dyer formula is proved here.}
The conditional counting implication is elementary; no novelty claim is
made for it.
\end{abstract}

\section{Scope and main results}\label{sec:main}

For an integer $n\ge5$ put
\begin{equation}\label{eq:params}
 a_n=\frac{n^2+27}{27-n^2},\qquad
 b_n=\frac{6n}{27-n^2},\qquad
 E_n:\ y^2=x^3-a_n^2x+b_n^2.
\end{equation}
Our parameter set is
\begin{equation}\label{eq:progression}
 n=5+759m,\qquad m=0,1,2,\ldots.
\end{equation}
The three rational points are
\begin{equation}\label{eq:points}
 P_n=(0,b_n),\quad Q_n=(a_n,b_n),\quad R_n=(-1,2b_n).
\end{equation}

\begin{theorem}[Three independent rational sections]\label{thm:three}
For every $n$ in \eqref{eq:progression}, $E_n$ is an elliptic curve
with $E_n(\Q)[2]=0$. The three classes
$[P_n],[Q_n],[R_n]\in E_n(\Q)/2E_n(\Q)$ are linearly independent
over $\F_2$. Consequently,
\[
 \rank E_n(\Q)\ge3,\qquad
 \dim_{\F_2}\Sel_2(E_n/\Q)\ge3.
\]
The rational function $j(E_n)$ is nonconstant; in particular, the
progression gives infinitely many distinct $\Q$-isomorphism classes
(and infinitely many distinct geometric $j$-invariants).
\end{theorem}

\begin{proposition}[An exceptional rank-at-least-four fiber]\label{prop:four}
At $n=5$ the curve and four rational points are
\[
 E_5:\ y^2=x^3-676x+225,
\quad P=(0,15),\ Q=(26,15),\ R=(-1,30),\ T=(-18,81).
\]
Their classes are linearly independent in $E_5(\Q)/2E_5(\Q)$.
Hence $\rank E_5(\Q)\ge4$ and
$\dim_{\F_2}\Sel_2(E_5/\Q)\ge4$.
\end{proposition}

Define, for positive integers $X$,
\begin{equation}\label{eq:averages}
 A(X)=\frac1X\sum_{m=0}^{X-1}
 \#\Sel_2(E_{5+759m}/\Q),
\quad
 N_3(X)=\#\{0\le m<X:\dim_{\F_2}\Sel_2(E_{5+759m}/\Q)=3\}.
\end{equation}

\begin{theorem}[Conditional mean-size criterion]\label{thm:mean}
For every positive integer $X$,
\[
 N_3(X)\ \ge\ \frac{16-A(X)}8\,X.
\]
In particular, if there exist $\varepsilon>0$ and $X_0$ such that
$A(X)\le16-\varepsilon$ for all $X\ge X_0$, then
$N_3(X)\ge(\varepsilon/8)X$ for all $X\ge X_0$.
For every curve counted by $N_3(X)$,
\[
 \rank E_n(\Q)=3\quad\text{and}\quad \Sha(E_n/\Q)[2]=0.
\]
\end{theorem}

The hypothesis of Theorem~\ref{thm:mean} is \emph{not} proved,
and the statement is an implication, not a theorem asserting the
existence of infinitely many Selmer-dimension-three specializations.

\section{Proof of the rational-section results}

The identity
\begin{equation}\label{eq:pell}
 a_n^2-3b_n^2=1
\end{equation}
is immediate on inserting \eqref{eq:params}.
It gives $P_n,Q_n,R_n\in E_n(\Q)$:
for $Q_n$, the right side is $b_n^2$, and for $R_n$ it
is $-1+a_n^2+b_n^2=4b_n^2$.
The short-Weierstrass discriminant is
\begin{align}
 \Delta(E_n)
 &=16(4a_n^6-27b_n^4)\notag\\
 &=16(108b_n^6+81b_n^4+36b_n^2+4)>0.\label{eq:delta}
\end{align}
Thus every parameter in \eqref{eq:progression} gives a smooth
elliptic curve.

For $p\in\{3,11,23\}$, the denominator $27-n^2$ is a $p$-adic
unit throughout \eqref{eq:progression}, and
\[
 (a_n,b_n)\equiv(26,15)\pmod p.
\]
The discriminant of $E_5$ equals $19748739664$, which is prime
to these three primes. Thus the reductions are good and independent
of the parameter $m$.

Modulo $11$ the 2-division cubic is $x^3+6x+5$. A check of
the eleven residue classes gives no root. Since the curve has
good reduction at $11$, it follows that $E_n(\Q)[2]=0$.

To prove the independence of $P_n,Q_n,R_n$, consider a nonempty
binary sum that would lie in $2E_n(\Q)$. Its reduction at a good
prime must lie in $2E_n(\F_p)$.
At $p=3$, the group is
\[
 E_n(\F_3)=\{O,(0,0),(1,0),(2,0)\}
 \cong(\mathbb Z/2\mathbb Z)^2,
\]
so $2E_n(\F_3)=\{O\}$. The marked points reduce to
$\bar P_n=(0,0)$ and $\bar Q_n=\bar R_n=(2,0)$.
Among the seven nonempty binary combinations,
only $Q_n+R_n$ could therefore be a double over $\Q$.

Modulo $23$ the cubic factors as
\[
 x^3-26^2x+15^2\equiv(x-4)(x^2+4x+7).
\]
The group $E_n(\F_{23})$ has $20$ elements and its doubling
image is exactly
\begin{align*}
 2E_n(\F_{23})=\{&
 O,(4,0),(20,8),(20,15),(5,11),(5,12),\\
 &(10,10),(10,13),(22,7),(22,16)\}.
\end{align*}
The addition law gives $\overline{Q_n+R_n}=(2,10)$, which
is not in this image. The remaining possibility is excluded.
The three classes are independent modulo doubling.
As there is no rational 2-torsion, the Mordell--Weil theorem
implies $\rank E_n(\Q)\ge3$, and the Kummer exact sequence yields
the Selmer lower bound.

Finally, the $j$-invariant is
\[
 j(E_n)=1728\frac{4a_n^6}{4a_n^6-27b_n^4}.
\]
Since $b_n\ne0$ for $n\ge5$, this value is never $1728$,
while $b_n\to0$ and $a_n^2\to1$ as $n\to\infty$,
so $j(E_n)\to1728$. A sequence of rational numbers with this
limit, none equal to the limit, cannot take only finitely many
values. This proves nonisotriviality and the last assertion of
Theorem~\ref{thm:three}. \qed

\medskip
\noindent\emph{Proof of Proposition~\ref{prop:four}.}
Direct substitution confirms all four points in the statement.
The same good reductions at $3$ and $23$, together with good
reduction at $5$, give the following exact sieve of the $15$
nonzero binary sums; masks use bit weights $1,2,4,8$ for
$P,Q,R,T$, respectively:
\[
\begin{array}{c|c|l}
p & \#E_5(\F_p) & \text{nonzero masks still potentially doubled}\\ \hline
3 &4 & 6,9,15\\
5 &8 &6\\
23&20&\varnothing
\end{array}
\]
At $p=5$ the doubling image is $\{O,(0,0)\}$;
at $p=23$ it is the ten-element set listed above.
Thus no nonempty binary sum is a double in $E_5(\Q)$.
Since $E_5(\Q)[2]=0$, the free rank is at least four,
and the Kummer sequence gives the asserted Selmer lower bound.
\hfill$\square$

\section{The conditional arithmetic-statistics argument}

\noindent\emph{Proof of Theorem~\ref{thm:mean}.}
Theorem~\ref{thm:three} shows that all Selmer dimensions
in the progression are at least three. Thus their cardinalities
belong to $\{8,16,32,64,\ldots\}$. Precisely $N_3(X)$ of
the first $X$ groups have cardinality $8$; each of the
others has cardinality at least $16$. Therefore
\[
 XA(X)\ge 8N_3(X)+16(X-N_3(X))
       =16X-8N_3(X).
\]
Rearranging proves the claimed inequality.
If $\dim_{\F_2}\Sel_2(E_n/\Q)=3$, then the Kummer
exact sequence
\[
0\longrightarrow E_n(\Q)/2E_n(\Q)
\longrightarrow\Sel_2(E_n/\Q)
\longrightarrow\Sha(E_n/\Q)[2]\longrightarrow0
\]
and Theorem~\ref{thm:three} give rank exactly $3$
and $\Sha(E_n/\Q)[2]=0$. \hfill$\square$

\begin{remark}[A genuine gap, not a statistical estimate]
The condition $\limsup_{X\to\infty}A(X)<16$ is an
unproved hypothesis. The preceding counting inequality
does not produce it. Nor can an average-size theorem
for all elliptic curves ordered by height automatically
be restricted to this thin arithmetic progression.
\end{remark}

\section{Cubic-field descent formulation and open problem}

Set $d_n=n^2-27$ and $A_n=n^2+27$. Clearing denominators,
the change of coordinates $U=d_n x$, $Z=d_n^2 y$ turns
$E_n$ into
\begin{equation}\label{eq:cubic}
 Z^2=d_n\,h_n(U),\qquad
 h_n(U)=U^3-A_n^2 U+36n^2d_n.
\end{equation}
The discriminant of the monic cubic $h_n$ is
\[
 \operatorname{disc}(h_n)
 =4\big((n^2+27)^6-8748n^4(n^2-27)^2\big).
\]
The irreducibility of the 2-division cubic modulo $11$
shows that
$K_n=\Q[U]/(h_n(U))=\Q(u_n)$, where $h_n(u_n)=0$, is a cubic number field.
In the usual cubic \'etale-algebra model for 2-descent,
the images of the marked points have representatives
\[
 [-d_nu_n],\quad
 [-d_n(A_n+u_n)],\quad
 [-d_n(d_n+u_n)]
 \quad\text{in }K_n^\times/K_n^{\times2}.
\]
Their independence follows from Theorem~\ref{thm:three};
it does not imply that they exhaust the Selmer group.

For each $n$, full 2-descent would require determination of
the admissibly ramified norm-one squareclasses, including
the varying cubic-field class-group 2-torsion, and the
local Kummer images at every relevant place
\cite{Schaefer2003}. Only after these are explicitly
constructed can one speak of an exact local-condition
matrix $\lambda_n$ whose kernel realizes $\Sel_2(E_n/\Q)$.

\begin{quote}
\textbf{Open problem (not a claim of this note).}
Prove that an explicitly specified infinite subset of
$n=5+759m$ has
$\dim_{\F_2}\Sel_2(E_n/\Q)=3$.
A sufficient arithmetic-statistics hypothesis is
$\limsup_{X\to\infty}A(X)<16$.
\end{quote}

The known general mean-size-$3$ result for 2-Selmer
groups of elliptic curves ordered by height
\cite{BhargavaShankar2015} does not provide this
thin-family estimate. Results on families with marked
points \cite{BhargavaHo2022} and rank nonjump
\cite{Zywina2025}, as well as the more recent
prescribed-rank theorem \cite{Zywina2026}, supply
important context but do not establish the above
Selmer exhaustion in this family.

\section{Reproducibility, independence and submission boundaries}

The exact finite checks underlying
Theorem~\ref{thm:three} and Proposition~\ref{prop:four}
are tracked in the MathOS source tree by
\texttt{mathos/research/tect\_x\_rank3\_family\_certificate.py}
and
\texttt{mathos/research/tect\_x\_rank4\_fiber\_certificate.py},
with adversarial regression checks in
\texttt{tests/test\_tect\_x\_noni\_selmer.py}.
These tests verify finite-field computations; they do
\emph{not} replace proofs of global arithmetic results.

The present document is a \emph{preliminary research
manuscript}, not a completed proof of infinitude
of exact-rank-three fibers or of Birch--Swinnerton-Dyer.
Independent specialist proof review, a full literature
and novelty audit, and authorship clearance are
pending. Its elementary conditional mean-size
criterion is not represented as a new deep theorem.

\begin{thebibliography}{9}
\bibitem{Schaefer2003}
E.~F.~Schaefer and M.~Stoll, \emph{How to do a $p$-descent on an elliptic curve},
Trans.\ Amer.\ Math.\ Soc.\ \textbf{356} (2004), 1209--1231.
\url{https://eref.uni-bayreuth.de/id/eprint/6228/}.

\bibitem{BhargavaShankar2015}
M.~Bhargava and A.~Shankar,
\emph{Binary quartic forms having bounded invariants, and the
boundedness of the average rank of elliptic curves},
Ann.\ of Math.\ \textbf{181} (2015), 191--242.
\url{https://annals.math.princeton.edu/2015/181-1/p03}.

\bibitem{BhargavaHo2022}
M.~Bhargava and W.~Ho,
\emph{On average sizes of Selmer groups and ranks in
families of elliptic curves having marked points},
\url{https://arxiv.org/abs/2207.03309}
(preprint; check current bibliographic status before submission).

\bibitem{Zywina2025}
D.~Zywina,
\emph{An elliptic surface with infinitely many fibers
for which the rank does not jump},
\url{https://arxiv.org/abs/2502.01026} (2025).

\bibitem{Zywina2026}
D.~Zywina, \emph{On the infinitude of elliptic curves
over a number field with prescribed small rank},
\url{https://arxiv.org/abs/2602.10865} (2026).
\end{thebibliography}

\end{document}
